\section{System Architecture}

The SkillChain platform follows a multi-layer architecture that separates concerns and enables scalability, maintainability, and security.

\subsection{High-Level Architecture}

\begin{table}[H]
\centering
\caption{System Architecture Layers}
\begin{tabular}{|p{3cm}|p{5cm}|p{5cm}|}
\hline
\textbf{Layer} & \textbf{Components} & \textbf{Responsibilities} \\
\hline
Presentation Layer & Next.js Frontend, Client UI, Freelancer UI, Dashboard & User interface rendering and interaction \\
\hline
Application Layer & Wagmi hooks, Zustand state, React Query, IPFS integration & Business logic and data management \\
\hline
Blockchain Layer & Ethereum Network, Smart Contracts & Decentralized computation and storage \\
\hline
Storage Layer & IPFS/Pinata & Decentralized content storage \\
\hline
\end{tabular}
\end{table}

\subsection{Layer 1: Presentation Layer}

The presentation layer is built using Next.js 16 with the App Router architecture. It provides:

\begin{itemize}
    \item Responsive user interfaces for clients and freelancers.
    \item Role-based dashboard views.
    \item Wallet connection management via RainbowKit.
    \item Real-time transaction status updates.
    \item Form validation and error handling.
\end{itemize}

\subsection{Layer 2: Application Layer}

The application layer handles business logic and state management:

\begin{itemize}
    \item \textbf{Wagmi Hooks:} Ethereum blockchain interactions (read/write operations).
    \item \textbf{Zustand:} Global state management for user sessions and application state.
    \item \textbf{React Query:} Server state management with caching and synchronization.
    \item \textbf{IPFS Integration:} Content upload and retrieval via Pinata API.
\end{itemize}

\subsection{Layer 3: Blockchain Layer}

The blockchain layer consists of three interconnected smart contracts:

\begin{enumerate}
    \item \textbf{UserProfile Contract:} Manages user registration and profiles.
    \item \textbf{JobsContract:} Handles job creation, escrow, and payment release.
    \item \textbf{ProposalsContract:} Manages proposal submission and approval.
\end{enumerate}

\subsection{Layer 4: Storage Layer}

IPFS provides decentralized storage for:

\begin{itemize}
    \item User profile metadata.
    \item Job descriptions and requirements.
    \item Proposal details and pricing.
    \item Work deliverables and attachments.
\end{itemize}

\section{Data Flow Diagrams}

\subsection{Job Creation Flow}

\begin{enumerate}
    \item Client fills job creation form with title, description, budget.
    \item Frontend uploads metadata to IPFS and receives CID.
    \item Client confirms transaction with ETH deposit.
    \item Smart contract creates job with IPFS CID and locks funds.
    \item Job is indexed and displayed in job listings.
\end{enumerate}

\subsection{Proposal and Hiring Flow}

\begin{enumerate}
    \item Freelancer views available jobs.
    \item Freelancer submits proposal with cover letter and bid.
    \item Proposal metadata uploaded to IPFS.
    \item Smart contract stores proposal reference.
    \item Client reviews proposals and approves selected freelancer.
    \item Approval automatically triggers freelancer hiring.
    \item Job status changes to ``In Progress''.
\end{enumerate}

\subsection{Completion and Payment Flow}

\begin{enumerate}
    \item Freelancer completes work and uploads deliverables to IPFS.
    \item Freelancer marks job complete via smart contract.
    \item Client reviews deliverables.
    \item Client marks job complete via smart contract.
    \item Upon mutual approval, funds are released to freelancer.
    \item Job marked as completed.
\end{enumerate}

\section{Entity Relationship Diagrams}

\subsection{Core Entities}

The system consists of three primary entities:

\begin{table}[H]
\centering
\caption{Entity Relationships}
\begin{tabular}{|p{3cm}|p{4cm}|p{5cm}|}
\hline
\textbf{Entity} & \textbf{Attributes} & \textbf{Relationships} \\
\hline
UserProfile & userID, username, cid, role, userAddress & One-to-Many with Jobs (as Client), One-to-Many with Proposals (as Freelancer) \\
\hline
Job & jobID, jobCID, deliverableCID, completed, client, freelancer, amount, clientApproved, freelancerApproved & Many-to-One with Client, Many-to-One with Freelancer, One-to-Many with Proposals \\
\hline
Proposal & proposalID, proposalCID, jobID, approved, freelancer & Many-to-One with Job, Many-to-One with Freelancer \\
\hline
\end{tabular}
\end{table}

\subsection{Data Relationships}

\begin{itemize}
    \item A \textbf{User} can be either a Client or Freelancer (Role-based).
    \item A \textbf{Client} can create multiple \textbf{Jobs} (One-to-Many).
    \item A \textbf{Job} can have multiple \textbf{Proposals} (One-to-Many).
    \item A \textbf{Freelancer} can submit multiple \textbf{Proposals} (One-to-Many).
    \item A \textbf{Job} is assigned to exactly one \textbf{Freelancer} (Many-to-One).
    \item Each \textbf{Job} has exactly one \textbf{Client} (Many-to-One).
\end{itemize}

\section{UML Diagrams}

\subsection{Use Case Diagram}

\textbf{Actors:}
\begin{itemize}
    \item Client
    \item Freelancer
    \item System (Smart Contracts)
\end{itemize}

\textbf{Client Use Cases:}
\begin{itemize}
    \item Register/Create Profile
    \item Create Job
    \item View Posted Jobs
    \item Review Proposals
    \item Approve Proposal/Hire Freelancer
    \item Mark Job as Completed
\end{itemize}

\textbf{Freelancer Use Cases:}
\begin{itemize}
    \item Register/Create Profile
    \item Browse Jobs
    \item Submit Proposal
    \item View Submitted Proposals
    \item Submit Deliverables
    \item Mark Job as Completed
\end{itemize}

\subsection{Class Diagram}

\textbf{Smart Contract Classes:}

\begin{lstlisting}[language=Solidity, caption={UserProfile Contract Structure}]
contract UserProfile {
    struct Profile {
        uint256 userID;
        string username;
        string cid;
        Role role;
        address userAddress;
    }
    
    function setProfile(...) external
    function getProfileByAddress(...) external view
    function getUserByUsername(...) external view
}
\end{lstlisting}

\begin{lstlisting}[language=Solidity, caption={JobsContract Structure}]
contract JobsContract {
    struct Job {
        uint256 jobID;
        string jobCID;
        string deliverableCID;
        bool completed;
        address client;
        address freelancer;
        uint256 amount;
        bool clientApproved;
        bool freelancerApproved;
    }
    
    function CreateJob(...) external payable
    function HireFreelancer(...) external
    function MarkFLJobCompleted(...) external
    function MarkClientJobCompleted(...) external
}
\end{lstlisting}

\begin{lstlisting}[language=Solidity, caption={ProposalsContract Structure}]
contract ProposalsContract {
    struct Proposal {
        uint256 proposalID;
        string proposalCID;
        uint256 jobID;
        bool approved;
        address freelancer;
    }
    
    function CreateProposal(...) external
    function ApproveProposal(...) external
    function getProposalsByJobID(...) external view
}
\end{lstlisting}

\subsection{Sequence Diagram: Job Creation}

\begin{enumerate}
    \item Client opens job creation form.
    \item Client fills form with job details.
    \item System validates form inputs.
    \item System uploads metadata to IPFS.
    \item IPFS returns Content Identifier (CID).
    \item Client confirms job creation with ETH deposit.
    \item MetaMask prompts for transaction confirmation.
    \item Smart contract receives transaction.
    \item Contract validates ETH amount.
    \item Contract creates job with CID and locks funds.
    \item Contract emits JobCreated event.
    \item Frontend updates UI with success message.
\end{enumerate}

\subsection{Activity Diagram: Complete Job Lifecycle}

The complete job lifecycle involves seven phases:

\textbf{Phase 1: Setup}
\begin{itemize}
    \item User A registers as Client.
    \item User B registers as Freelancer.
\end{itemize}

\textbf{Phase 2: Job Creation}
\begin{itemize}
    \item Client creates job with ETH escrow.
    \item Contract stores job and emits event.
\end{itemize}

\textbf{Phase 3: Proposal Submission}
\begin{itemize}
    \item Freelancer submits proposal.
    \item Contract stores proposal reference.
\end{itemize}

\textbf{Phase 4: Proposal Approval}
\begin{itemize}
    \item Client reviews and approves proposal.
    \item Contract automatically hires freelancer.
\end{itemize}

\textbf{Phase 5: Work Execution}
\begin{itemize}
    \item Freelancer completes work off-chain.
\end{itemize}

\textbf{Phase 6: Completion Marking}
\begin{itemize}
    \item Freelancer marks job complete.
    \item Client marks job complete.
\end{itemize}

\textbf{Phase 7: Payment Release}
\begin{itemize}
    \item Upon mutual approval, funds released.
    \item Job marked as completed.
\end{itemize}

\subsection{State Diagram: Job States}

\begin{table}[H]
\centering
\caption{Job State Transitions}
\begin{tabular}{|p{4cm}|p{4cm}|p{5cm}|}
\hline
\textbf{Current State} & \textbf{Event} & \textbf{Next State} \\
\hline
Initial & CreateJob transaction & Open \\
\hline
Open & ApproveProposal & In Progress \\
\hline
In Progress & MarkFLJobCompleted & Awaiting Client Approval \\
\hline
In Progress & MarkClientJobCompleted & Awaiting Freelancer Approval \\
\hline
Awaiting Client Approval & MarkClientJobCompleted & Completed \\
\hline
Awaiting Freelancer Approval & MarkFLJobCompleted & Completed \\
\hline
\end{tabular}
\end{table}